\begin{table}[tbp]
  \centering
  \small
  \caption{v3 配对响应与 v4 阶段验证的历史证据}
  \label{tab:v3-v4-historical-evidence}
  \begin{tabularx}{\linewidth}{@{}l>{\raggedright\arraybackslash}X@{}}
    \toprule
    证据项 & 结果与适用范围 \\
    \midrule
    v3 配对响应 &
    64 个独立事件产生 128 条 bit0/bit1 配对响应，分类按响应计数。
    在 5\,mV 主阈值下，96 条正确、20 条判据误读、8 条临界、4 条数值不确定。
    结果限于未校准二维代表模型的冻结 5\,MeV 工况。 \\
    v4 核机制与三维源 &
    厚硅诊断的 212 个初级事件中记录到 9 个合格的质子非弹性反应见证。
    3 个三维源账本完成离线闭合，器件注入资格仍待验证。 \\
    v4 细网格通过项 &
    细网格 A0/A1 源检查、A1 全部 205 帧的源场测量、
    以及各用同网格 A0 作对照的粗细网格比较通过。
    该资格限于本次二维源与连续响应检查。 \\
    v4 保留的空间失败 &
    粗网格 A1 的中心窗口生成量占比绝对误差约 0.0595，
    超过 0.05 门槛，原失败记录保留。 \\
    v4 保留的旧失败 &
    旧时间步敏感性检查中，bit1 存储节点与位线差异约为
    1.709\,mV 和 5.140\,mV，均超出 1\,mV 容差；端口总电荷差约为
    1.630\%，低于 2\% 容差，整项仍判为失败。
    旧源输出记录有 322 次写入，最终仅保留 2 个 TDR 文件。
    旧 \(\alpha\) 试验的最大接受步约为 56.17\,fs，超过 40\,fs 约束。 \\
    \bottomrule
  \end{tabularx}
  \par\smallskip
  \begin{minipage}{\linewidth}
    \footnotesize
    证据日期（UTC）为 v4 阶段收口 2026-10-13、历史复用审计 2026-11-16。
    来源为共享证据索引、《现有成果的复用证据》及《DRAM SEE v4 阶段执行报告》。
    v3 分类适用于原读出判据。v4 证据覆盖机制观测与局部接口检查，
    器件校准、三维器件响应、翻转截面与训练资格仍待验证。
  \end{minipage}
\end{table}
